% !TITLE 楚辞·天问（节选）
% !AUTHOR 屈原
% !DYNASTY 先秦
% !GENRE 楚辞
% !ORDER 12
% !SOURCE 楚辞·天问

\begin{poem}
曰：遂古之初\textsuperscript{1}，谁传道之\textsuperscript{2}？
上下未形\textsuperscript{3}，何由考之\textsuperscript{4}？
冥昭瞢暗\textsuperscript{5}，谁能极之\textsuperscript{6}？
冯翼惟象\textsuperscript{7}，何以识之？
明明暗暗\textsuperscript{8}，惟时何为\textsuperscript{9}？
阴阳三合\textsuperscript{10}，何本何化\textsuperscript{11}？
圜则九重\textsuperscript{12}，孰营度之\textsuperscript{13}？
惟兹何功\textsuperscript{14}，孰初作之？
斡维焉系\textsuperscript{15}？天极焉加\textsuperscript{16}？
八柱何当\textsuperscript{17}？东南何亏\textsuperscript{18}？
九天之际\textsuperscript{19}，安放安属\textsuperscript{20}？
隅隈多有\textsuperscript{21}，谁知其数？
天何所沓\textsuperscript{22}？十二焉分\textsuperscript{23}？
日月安属\textsuperscript{24}？列星安陈\textsuperscript{25}？
\end{poem}

\begin{annotations}
\notetext{1}{遂古之初：远古的开端。遂，通"邃"，深远。一切开始的那个时刻。}
\notetext{2}{谁传道之：谁将此事传说下来？传道，传说、叙述。天地初开时还没有人，那后来的记载从何而来？}
\notetext{3}{上下未形：天地还没有成形。上指天，下指地。}
\notetext{4}{何由考之：凭什么去考究它？考，考察、探究。天地混沌不分时，谁也无法观察。}
\notetext{5}{冥昭瞢暗：白天黑夜还是一片昏暗。冥，黑夜；昭，白昼；瞢（méng）暗，模糊不清。}
\notetext{6}{谁能极之：谁能穷究它的真相？极，穷尽。}
\notetext{7}{冯翼惟象：元气满溢充盈，只有无形无状的"象"。冯（píng），充满；翼，弥漫。古人认为宇宙初开时只有混沌充溢的元气，无形无质。}
\notetext{8}{明明暗暗：光明与黑暗交替。}
\notetext{9}{惟时何为：这是因为什么？昼夜交替的规律是怎么形成的？}
\notetext{10}{阴阳三合：阴阳二气参合。三，通"参"，参杂、相合。}
\notetext{11}{何本何化：哪个是本原，哪个是化生出来的？阴和阳，哪个最先存在，哪个由对方化生？}
\notetext{12}{圜则九重：圆形的天盖有九重。圜，同"圆"，指天。古人认为天是圆的，有九层。}
\notetext{13}{孰营度之：是谁营造并度量了它？营，建造；度（duó），测量。}
\notetext{14}{惟兹何功：这是何等巨大的工程！功，工程。}
\notetext{15}{斡维焉系：旋转的枢纽系在哪里？斡（wò），旋转；维，大绳。古人认为天旋转不息，那转轴又拴在什么东西上？}
\notetext{16}{天极焉加：天的极顶加在哪里？天极，天的最高处；加，安放。}
\notetext{17}{八柱何当：八根擎天柱立在什么地方？古人认为天由八根巨柱支撑。当，在、位于。}
\notetext{18}{东南何亏：东南方为什么低陷了？亏，亏损、低陷。中国地势西北高东南低，古人用地陷来解释。}
\notetext{19}{九天之际：九重天的边界。际，边际。}
\notetext{20}{安放安属：安置在哪里、连接到哪里？属（zhǔ），连接。}
\notetext{21}{隅隈多有：天的角落和弯曲处很多。隅（yú），角落；隈（wēi），山水弯曲处。}
\notetext{22}{天何所沓：天在哪里与地会合？沓（tà），汇合。古人认为天地在远方相接。}
\notetext{23}{十二焉分：十二辰又是怎么划分的？十二，指黄道十二辰（太阳在天空中运行的十二个位置）。}
\notetext{24}{日月安属：日月附着在哪里？属（zhǔ），附着。日月悬在空中，它们挂靠在哪里？}
\notetext{25}{列星安陈：群星是怎样陈列的？陈，排列、陈列。为什么星星会固定在各自的位置上？}
\end{annotations}

\begin{appreciation}
《天问》是屈原最奇特的作品，也是中国诗歌史上独一无二的存在。全诗三百七十四句，提出了一百七十多个问题，从宇宙的起源问到人间的兴亡，从神话的真伪问到历史的公正。没有一句回答，全是追问。这里节选的开篇部分，问的是宇宙的开端——在两千三百年前，屈原已经把目光投向了天地初开的那一瞬间。

曰：遂古之初，谁传道之？开篇便是惊心动魄。远古的尽头，天地还没有分开，一片混沌。屈原问：那是谁把这一切告诉后来人的？这个问题直指人类认知的边界：既然没有人亲眼见过宇宙的诞生，那所有关于创世的叙述，依据何在？这是一个哲学家的追问，也是一个诗人的追问——它不为求答案，只为让人意识到自己的无知和渺小。

接着是一连串关于天地构造的问题。圜则九重，孰营度之？天有九重，是谁测度的？斡维焉系？天在旋转，那根转轴拴在哪里？八柱何当？八根擎天柱在什么地方立着？东南何亏？东南地势为什么塌陷了？日月安属？列星安陈？日月星辰又是怎么挂在天上的？这些问题在今天看来是科学问题，但屈原提出它们的方式是完全诗性的。他没有试图测量、计算、推演，他只是问——像一个孩子看着星空时提出的问题。但这种孩子式的发问恰恰是最深刻的，因为它回到了人类认知的起点：在一切知识建立之前，我们首先要问一个问题。

《天问》的伟大，不在于它提出了什么高明的答案，而在于它敢于提出这些没有答案的问题。在那个巫师和祭司解释一切的时代，屈原说：我不相信那些现成的答案。我要自己来问。这种质疑一切的精神，比任何答案都珍贵。两千年后，当现代科学追问宇宙起源时，它所依赖的，不也是这种"敢问"的精神吗？

有意思的是，《天问》全篇没有一句屈原对自己的描写，但他的形象反而比任何自述都更鲜明——我们看到一个孤独的人，仰天而问，把问题一个一个地掷向苍穹，不期待回音，只因为不能不问。这便是人类精神最辉煌的姿态：面对无限的宇宙，我们不放弃追问的权利。
\end{appreciation}
